% 第 3 章：严格多时段 LinDistFlow 恢复 MILP
\section{严格多时段 LinDistFlow 恢复 MILP}\label{sec:res-mip}

\subsection{功能定位}
\fld{MultiPeriodMIPLinDistFlow}（\srcpath{run\_distribution\_resilience\_mip\_assessment}，
\srcpath{resilience\_assessment.hpp:458}）在整个时域上构建\textbf{二元} MILP——支路开关变量、
母线带电变量、单商品流连通与 LinDistFlow 电压——用注册 MILP 后端
（Native B\&C / HiGHS / Gurobi，\fld{DistributionResilienceMIPSolver}，:22）求解，\textbf{不}回退到
启发式路径。选项见
\texttt{DistributionResilience\allowbreak MIPOptions}（:36）。

\subsection{决策变量}
对支路 $e$、母线 $i$、时步 $t$：闭合状态 $\beta_{e,t}\in\{0,1\}$、母线带电 $z_{i,t}\in\{0,1\}$、
切负荷 $s_{i,t}\ge0$、支路有功/无功 $P_{e,t},Q_{e,t}$、平方电压 $v_{i,t}=|V_{i,t}|^2$、
储能 SOC 与 MESS 路由/出力变量。

\subsection{约束}
\begin{align}
&\text{支路潮流限：} &|P_{e,t}|&\le \bar S_e\,\beta_{e,t},\quad \bar S_e=\text{\fld{default\_branch\_rate\_mva}},\\
&\text{LinDistFlow 压降：} &\big|v_{j,t}-v_{i,t}+2(r_eP_{e,t}+x_eQ_{e,t})\big|&\le M(1-\beta_{e,t}),\ M=\text{\fld{big\_m\_voltage}},\\
&\text{电压包络：} &\underline v^2 z_{i,t}\le v_{i,t}&\le \overline v^2 z_{i,t},\ [\underline v,\overline v]=[\text{\fld{v\_min\_pu}},\text{\fld{v\_max\_pu}}],\\
&\text{节点平衡：} &\textstyle\sum_{e\in\delta(i)}\!\pm P_{e,t}+P_{i,t}^{src}&=D_{i,t}-s_{i,t},\\
&\text{切负荷界：} &0\le s_{i,t}&\le D_{i,t},\\
&\text{连通/辐射：} &\text{单商品流（Dantzig）}&+\text{辐射森林}.
\end{align}
辐射性允许多根森林（\fld{allow\_multi\_root\_forest}）与构网储能/MESS 作根
（\fld{allow\_grid\_forming\_storage\_roots}、\fld{allow\_mess\_black\_start}）。DC 侧、VSC/DC-DC 传输
按有效性标志纳入（默认 DC 支路流限与换流器传输限\textbf{未}强制，见第~\ref{sec:res-overview}~章）。

\subsection{目标函数}
最小化按优先级加权的失供加操作代价：
\begin{equation}
\min\ \sum_t\Big[\sum_i \pi(i)\,s_{i,t}
+c_{sw}\!\sum_e|\Delta\beta_{e,t}|
+c_{mess}\,\mathrm{km}_t
+c_{curt}\,P^{curt}_t\Big],
\end{equation}
其中优先级罚权 $\pi(i)\in\{$\fld{shed\_penalty\_critical}$=10^6$, \fld{shed\_penalty\_high}$=10^4$,
\fld{shed\_penalty\_medium}$=10^2$, \fld{shed\_penalty\_low}$=1\}$，
$c_{sw}=$\fld{switching\_cost}、$c_{mess}=$\fld{mess\_travel\_cost\_per\_km}、
$c_{curt}=$\fld{renewable\_curtailment\_cost}。

\subsection{求解与规模}
Native B\&C 支持 \fld{native\_node\_selection}（Hybrid/BestFirst）与可选 CGLP 提升-投影析取割
（\fld{enable\_cglp\_cuts}）；停止判据为 \fld{mip\_gap}（默认 0.02）、\fld{max\_time\_s}（120）、
\fld{max\_nodes}（50000）。模型规模与状态记入 \fld{DistributionResilienceModelStats}（:69）：
变量/二元/整数数、约束数、\fld{objective\_value}、\fld{mip\_gap}、\fld{solver\_status}、
\fld{cglp\_cuts\_added}。\fld{mip\_gap\_within\_tolerance} 为近似最优证书，非全局最优。
